\documentclass[11pt,a4paper]{article}

% -------------------------------------------------------------------------

% 한국어 설정
\usepackage{kotex}

% 링크 설정
\usepackage[hidelinks,colorlinks=true,linkcolor=magenta,citecolor=magenta,urlcolor=blue]{hyperref}

% 그림 설정
\usepackage{graphicx}
\usepackage{caption}
\usepackage{subcaption}
\setlength{\textfloatsep}{25pt plus 2pt minus 4pt}
\setlength{\floatsep}{25pt plus 2pt minus 2pt}
\setlength{\intextsep}{25pt plus 2pt minus 2pt}

% 표 설정
\usepackage{booktabs}

% 수식 설정
\usepackage{amsmath}
\usepackage{amssymb}
\usepackage{mathtools}

% 여백 설정
\usepackage{geometry}
\geometry{margin=25mm}

% 섹션 설정
\usepackage{titlesec}
\titlespacing*{\section}{0pt}{5ex plus 1ex minus .2ex}{2ex}
\titlespacing*{\subsection}{0pt}{4ex plus 1ex minus .2ex}{1ex}
\titlespacing*{\subsubsection}{0pt}{3ex plus .8ex minus .2ex}{1ex}
\setcounter{secnumdepth}{3}
\titleformat{\section}{\large\bfseries}{\thesection}{0.5em}{}
\titleformat{\subsection}{\normalsize\bfseries}{\thesubsection}{0.5em}{}
\titleformat{\subsubsection}{\normalsize\bfseries}{\thesubsubsection}{0.5em}{}

% 행간 설정
\usepackage{setspace}
\setstretch{1.2}

% 들여쓰기 설정
\usepackage{indentfirst}

% -------------------------------------------------------------------------

% 제목 등 설정
\makeatletter
\newcommand{\subtitle}[1]{\gdef\@subtitle{#1}}
\newcommand{\studentinfo}[1]{\gdef\@studentinfo{#1}}
\def\@subtitle{}
\def\@studentinfo{}
\renewcommand{\maketitle}{%
  \begin{center}
    {\LARGE\bfseries \@title\par}
    \vspace{0.5em}
    {\large \@subtitle\par}
  \end{center}
  \vspace{-1.0em}
  \begin{flushright}
    \normalsize \@studentinfo
  \end{flushright}
}
\makeatother

% -------------------------------------------------------------------------
% -------------------------------------------------------------------------

\begin{document}

\title{과제1 보고서}

\subtitle{}

\studentinfo{학번 이름}

\maketitle

% -------------------------------------------------------------------------

\section{서론}

    과제 1은 대부분의 함수를 기존에 알고 있던 개념을 바탕으로 어렵지 않게 구현할 수 있었다.
    다만 \texttt{is\_prime()}과 \texttt{generate\_primes()} 두 함수를 최적화하는 과정에서
    처음 접하는 개념이 등장하여 추가적인 이해가 필요했다.

% -------------------------------------------------------------------------

\section{\texttt{is\_prime()}}

    \subsection{초기 구현 및 문제 발생}

        처음에는 소수 판별을 위해 2부터 \texttt{value}까지 모든 수로 나누는 방식을 구현했다.
        그러나 \texttt{value}가 최대 10,000,000일 경우 반복 횟수가 지나치게 많아져 성능 문제가 발생했다.
        이후 예외 케이스를 사전에 처리하고 홀수만 순회하도록 개선했지만, 여전히 최적의 방법은 아니라고 판단했다. 

    \subsection{이해 및 문제 해결}

        $\texttt{value} = a \times b$ 일 때, $a$와 $b$ 중 하나는 반드시
        $\sqrt{\texttt{value}}$ 이하이고, 그렇기에 범위가 $\sqrt{\texttt{value}}$까지여도 소수 여부를 판별할 수 있다는 개념을 알게 되었다.
        이를 바탕으로 반복문의 범위를 \texttt{int(value**0.5)+1} 까지로 설정해 불필요한 반복을 줄여 시간초과 문제를 
        해결할 수 있었다.

% -------------------------------------------------------------------------

\section{\texttt{generate\_primes()}}

    \subsection{초기 구현 및 문제 발생}

        처음에는 \texttt{is\_prime()}을 반복 호출하는 방식으로 단순하게 구현했다.
        \texttt{is\_prime()}을 최적화해 두었고, 이를 사용했기에 문제가 없을 것이라 생각했지만,
        실제로는 시간 초과(Time Limit Exceeded)가 발생했다.

    \subsection{이해 및 문제 해결}

        더 빠른 방법을 찾는 과정에서 '에라토스테네스의 체'를 처음 접했다.
        구현 중 배열을 \texttt{True}로 전체 초기화한다는 것이 낯설었는데, 제거되지 않은 수를 소수로 가정해야 판별이 가능함을 알고 나니 구현 원리가 직관적으로 느껴졌다.
        또 배수 제거를 $i \times i$부터 시작하는 부분도 바로 이해되지 않았는데,
        $i \times i$ 미만의 배수들은 이미 그보다 작은 소수들이 처리했기 때문에, 누락 없이 모든 합성수를 제거할 수 있다는 것을 깨닫고 받아들일 수 있었다.

% -------------------------------------------------------------------------

\section{결론}

        이번 과제를 통해 코드가 잘 동작하는 것에서 끝이 아니라, 효율적인 알고리즘을 고민하는 경험을 할 수 있었다.
        특히 수학적 성질을 코드에 직접 적용하거나 전혀 새로운 방식으로 문제를 접근해야 하는 것이 새롭고 어색했지만, 
        알고리즘 설계의 중요성을 체감할 수 있어서 유익했던 것 같다.
        
% % -------------------------------------------------------------------------

\end{document}
